\subsection{DRO improves worst-group accuracy under appropriate regularization}\label{sec:experiments_complexity_control}

Classically, we can control the generalization gap with regularization techniques that
constrain the model family's capacity to fit the training data.
In the modern overparameterized regime,
explicit regularization is not critical for average performance:
models can do well on average even when all regularization is removed \citep{zhang2017understanding},
and default regularization settings (like in the models trained above) still allow models to perfectly fit the training data.
Here, we study if increasing regularization strength---until the models no longer perfectly fit the training data---can rescue worst-case performance.
We find that departing from the vanishing-training-loss regime allows
DRO models to significantly outperform ERM models on worst-group test accuracy
while maintaining high average accuracy.
We investigate two types of regularization:

\paragraph{$\Ltwo$ penalties.}
The default coefficient of the $\Ltwo$-norm penalty $\lambda \|\theta\|_2^2$ in ResNet50 is $\lambda = 0.0001$ \citep{he2016resnet}.
We find that increasing $\lambda$ by several orders of magnitude---to $\lambda = 1.0$ for Waterbirds and $\lambda = 0.1$ for CelebA---does two things: 1) it prevents both ERM and DRO models from achieving perfect training accuracy, and 2) substantially reduces the generalization gap for each group.

With strong $\Ltwo$ penalties, both ERM and DRO models still achieve high average test accuracies.
However, because no model can achieve perfect training accuracy in this regime,
ERM models sacrifice worst-group training accuracy
($35.7\%$ and $40.4\%$ on Waterbirds and CelebA; \reftab{analysis}, \reffig{train}) and consequently obtain poor worst-group test accuracies ($21.3\%$ and $37.8\%$, respectively).

In contrast,
DRO models attain high worst-group training accuracy ($97.5\%$ and $93.4\%$ on Waterbirds and CelebA).
The small generalization gap in the strong-$\Ltwo$-penalty regime means that high worst-group training accuracy translates to high worst-group test accuracy,
which improves over ERM from $21.3\%$ to $84.6 \%$ on Waterbirds
and from $37.8\%$ to $86.7\%$ on CelebA.

While these results show that strong $\Ltwo$ penalties have a striking impact on ResNet50 models for Waterbirds and CelebA,
we found that increasing the $\Ltwo$ penalty on the BERT model for MultiNLI resulted in similar or worse robust accuracies than the default BERT model with no $\Ltwo$ penalty.

\paragraph{Early stopping.}
A different, implicit form of regularization is early stopping \citep{hardt2016train}.
We use the same settings in \refsec{experiments_unreg},
but only train each model for a fixed (small) number of epochs (\refsec{model_details}).
As with strong $\Ltwo$ penalties, curtailing training reduces the generalization gap
and prevents models from fitting the data perfectly.
In this setting, DRO also does substantially better than ERM
on worst-group test accuracy,
improving from
$6.7\%$ to $86.0\%$ on Waterbirds,
$25.0\%$ to $88.3\%$ on CelebA,
and $66.0\%$ to $77.7\%$ on MultiNLI.
Average test accuracies are comparably high in both ERM and DRO models,
though there is a small drop of $1-3\%$ for DRO (\reftab{analysis}, \reffig{train}).

\paragraph{Discussion.}
We conclude that regularization---preventing the model from perfectly fitting the training data---does matter for worst-group accuracy.
Specifically, it controls the generalization gap for each group, even on the worst-case group. Good worst-group test accuracy then becomes a question of good worst-group training accuracy.
Since no regularized model can perfectly fit the training data, ERM and DRO models make different training trade-offs:
ERM models sacrifice worst-group for average training accuracy and therefore have poor worst-group test accuracies,
while DRO models maintain high worst-group training accuracy and therefore do well at test time.
Our findings raise questions about the nature of generalization in neural networks,
which has been predominantly studied only in the context of average accuracy \citep{zhang2017understanding, hoffer2017train}.
